% Version 0.29. Actual modular support and the two-link construction.
% Pinned source: 7cbf7a0c32fca1badf05931879a37c3f8a2b03b4.
\section{Adjoint vanishing, modular support, and the two-link construction}
\label{sec:symplectic-modular-support}

\status{This chapter proves the literal \texttt{AdjointLiftingTheorem}
and the characteristic-zero fixed-weight support argument.  It also
constructs the two comparison links and their actual maps from universal
ordinary determinant rings.  Its final tower assertion retains the
precisely stated ordinary modular-quotient and characteristic-$\ell$
support inputs.  Those inputs are not consequences of adjoint Selmer
vanishing.  In particular no complete \texttt{UnramifiedInputs},
\texttt{HypothesesH'} or \texttt{ModUAlt} package is inferred from the
adjoint theorem.}

\subsection{Conventions and the independent adjoint input}

We retain the original algebraic closure $\Omega$, coefficient subrings
$O\subset\overline{\mathbb Q}_\ell$, and their specified inclusions.
For a CM field $L\subset\Omega$ put $L^+=\texttt{Lp}\,L$.
A source level $V$ gives a quotient
\[
 G_{L^+}\twoheadrightarrow\Gamma,\qquad
 G_L\twoheadrightarrow\Gamma_0=\ker(\eta),\qquad
 \Gamma'=\ker(\varepsilon|_{\Gamma_0}),
 \quad \eta:\Gamma\longrightarrow\{\pm1\}.
 \tag{QX.1}
\]
The maps commute with the $\ell$-adic cyclotomic character $\chi$.
The chosen complex conjugation $c\in\Gamma$ satisfies $c^2=1$ and
$\eta(c)=-1$.  For a representation $\rho:\Gamma_0\to
\operatorname{GL}_n(\overline{\mathbb Q}_\ell)$ the source's
\texttt{TypeI} is irreducibility of $\rho|_{\Gamma'}$, over the
algebraically closed coefficient field.  It is not a residual-image
hypothesis.

The source's \texttt{IsAutCsdA}$\,\iota\rho$ includes an actual regular
algebraic conjugate-self-dual cuspidal representation $\pi$ of
$\operatorname{GL}_n(\mathbb A_L)$ and its \texttt{AttC} attachment.
At an unramified place of norm $q$, this attachment compares the
characteristic polynomial of geometric Frobenius with
$q^{(n-1)/2}$ times the unitary Satake roots.  We keep this convention
throughout.  The extra \texttt{WitA} condition will matter in the
automorphic descent below, but is unnecessary in QX1.

We use Thorne, \emph{The vanishing of adjoint Bloch--Kato Selmer groups
of irreducible automorphic Galois representations}, Theorem 6.1,
\url{https://arxiv.org/abs/2207.04925}.  Its precise statement is:
for a CM field $L$, $n\ge2$, a regular algebraic cuspidal polarized
automorphic pair $(\pi,\xi)$, and an isomorphism
$\iota:\overline{\mathbb Q}_\ell\simeq\mathbb C$, irreducibility of
$r_{\pi,\iota}|_{G_{L(\mu_{\ell^\infty})}}$ implies
$H^1_f(L^+,\operatorname{ad}r_{\pi,\iota})=0$.
Here the adjoint is the polarized, or unitary, adjoint over $L^+$.
There is no condition on the residual image and no lower bound on
$\ell$.  The Galois-representation normalization is that of
Barnet-Lamb--Gee--Geraghty--Taylor, \emph{Potential automorphy and
change of weight}, Theorem 2.1.1,
\url{https://annals.math.princeton.edu/wp-content/uploads/annals-v179-n2-p03-p.pdf}.
Both are used with their hypotheses, rather than as an assertion that
all ordinary points are automorphic.

% BEGIN QX1
\begin{proposition}[QX1: the complete adjoint-lifting input]
\label{claim:QX1}
Let $n\ge2$, and let $\rho:\Gamma_0\to
\operatorname{GL}_n(\overline{\mathbb Q}_\ell)$ be continuous.
Suppose that $\rho\circ\texttt{pr}$ satisfies
\texttt{IsAutCsdA}$\,\iota$ and that $\rho|_{\Gamma'}$ is
irreducible.  Then the literal \texttt{AdZero}$\,V\rho$ holds.
Consequently the source's entire \texttt{AdjointLiftingTheorem}
is supplied, for every one of its original level, residual, point,
coefficient and embedding parameters.
\end{proposition}
\begin{proof}
First, the surjection $G_L\to\Gamma_0$ maps
$G_{L(\mu_{\ell^\infty})}$ \emph{onto} $\Gamma'$: lift an element
of $\Gamma'$, and use cyclotomic compatibility to see that its lift
has $\chi=1$.  Thus the irreducibility hypothesis is precisely the
one required by Thorne's theorem.  In particular $\rho$ and its
inflation to $G_L$ are absolutely irreducible.

Choose the actual $\pi$ in \texttt{IsAutCsdA}.  The quadratic
Hecke character $\delta_{L/L^+}$ is trivial after pullback by the
norm from $L$, so $\pi$, together with
$\delta_{L/L^+}^{\,n}$, is a polarized automorphic pair.  At each
real place this character has value $(-1)^n$ on $-1$, the required
parity.  The Galois multiplier is
\[
       \nu=\delta_{L/L^+}^{\,n}\chi^{\,1-n};
       \qquad \nu(c)=-1.
       \tag{QX.2}
\]
The displayed exponent and sign are also stated in Thorne's
\S6.3.  The \texttt{AttC} Frobenius convention above is exactly
the normalization of $r_{\pi,\iota}$.  Chebotarev and the
semisimple-character criterion identify $\rho\circ\texttt{pr}$
with $r_{\pi,\iota}$: their characteristic polynomials agree at
all but finitely many unramified places, hence on the whole compact
Galois group by continuity.  Both representations are semisimple.

Now take \emph{any} $P$ and $R$ occurring in \texttt{AdZero}.
The exact form and action identities are
\[
 \rho(g)^tP\rho(cgc^{-1})=\chi(g)^{1-n}P,
 \quad R(g)X=\rho(g)X\rho(g)^{-1},
 \quad R(c)X=-P^{-1}X^tP.
 \tag{QX.3}
\]
The first identity is the source's \texttt{IsForm}, and the latter
two are \texttt{IsAdExt}.  A standard polarization of
$r_{\pi,\iota}$ gives another intertwiner $P_\pi$ in the first
identity.  Absolute irreducibility and Schur's lemma give
$P=aP_\pi$ for a scalar $a\ne0$.  This scalar cancels in the last
formula.  Therefore, after inflation to $G_{L^+}$, the given $R$
is exactly the polarized adjoint in Thorne's theorem, including
the minus sign at complex conjugation.  No smaller symplectic Lie
algebra or unpolarized adjoint is substituted for it.

All these assertions can be made over one finite extension
$E/\mathbb Q_\ell$.  Here is the coefficient point explicitly.
The finite subextensions of $\overline{\mathbb Q}_\ell$ have a
countable cofinal family.  The inverse images of
$\operatorname{GL}_N(E_j)$ under a continuous representation are
closed subgroups covering the compact source group.  Baire gives
one with interior; adjoining the finitely many matrices of coset
representatives then captures the entire image over a finite
extension.  Apply this to the representation and, for a continuous
cocycle $z$, to its affine block representation.  Adjoin the entries
of $P$ and of the chosen intertwiner.  The action $R$ is continuous
on the two cosets of $\Gamma_0$ by (QX.3).  This is the same finite
coefficient capture used in RV4, with no completion of the
algebraic closure assumed.

The hypothesis \texttt{Af.IsBK} on the inflated $z$ is exactly
membership in $H^1_f(L^+,R)$.  Away from $\ell$ its stated local
primitive is the vanishing of the inertia restriction.  Above
$\ell$, KB6 identifies the original decomposition group with the
absolute Galois group of the completion, and BX1--BX4 identify
\texttt{finiteCocycles} with the kernel after tensoring with
$B_{\rm cris}$.  These comparisons commute with the original
restriction maps and with finite coefficient extension.  Descent
back to $E$ is harmless: a finite extension that contains any
needed primitive admits an $E$-linear projection taking $1$ to $1$;
equivalently, the local finite-condition maps commute with finite
flat coefficient extension.  Thus the inflated class belongs to
the exact Selmer group in Thorne's theorem and is zero.

Choose a matrix $X$ giving its coboundary.  Surjectivity
$G_{L^+}\to\Gamma$ gives
$z(g)=R(g)X-X$ for every $g\in\Gamma$, which is precisely
\texttt{AdZero}.  The literal \texttt{AdjointLiftingTheorem}
quantifies only over $2\le n\le4$; its point record already supplies
continuity, \texttt{IsAutCsdA} and \texttt{TypeI}.  Its additional
\texttt{Par}, \texttt{Free}, \texttt{N3}, residual and level
conditions need not be strengthened or used.  There are no rank-zero
or rank-one cases hidden in that literal statement.
\end{proof}
% END QX1

\subsection{Fixed-weight tangent vanishing and local support}

An ordinary determinant ring $B$ below is the actual
\texttt{UniversalOrdinaryDeterminantRing} of the source.  Its
universal datum consists of a degree-$n$ determinant law $D_B$ and
ordered local characters $\psi_{v,i}$.  It retains the residual
determinant, local residual characters and prescribed torsion.
Its \texttt{P} clause equates the characteristic coefficients of
each group element and its polarized star; its \texttt{Ord} and
\texttt{I} clauses are trace identities after left multiplication
by every group element.  These are the weaker universal
\texttt{Pair} relations of KX8.  They are not being replaced by
the full scalar-extended relations in \texttt{RingProps}.  A weight object
$\Lambda$ represents exactly the allowed restrictions of the
$\psi_{v,i}$ to inertia.  KX supplies
\[
       \Lambda\simeq O[[T_1,\ldots,T_N]],
       \qquad N=n[L^+:\mathbb Q].
       \tag{QX.4}
\]
Thus fixing a homomorphism from $\Lambda$ fixes all those ordered
inertial characters, not just their product.

% BEGIN QX2
\begin{lemma}[QX2: a complete commutative-algebra support criterion]
\label{claim:QX2}
Let $B$ be a complete Noetherian local $O$-algebra with the original
finite residue field, $f:\Lambda=O[[T_1,\ldots,T_N]]\to B$, and
$J\subset B$.  Assume that $B/J$ is finite and torsion-free as a
$\Lambda$-module.  Let $x:B\to\overline{\mathbb Q}_\ell$ be an
$O$-algebra map killing $J$.  If every lift
$x':B\to\overline{\mathbb Q}_\ell[\epsilon]/(\epsilon^2)$ of $x$
whose derivative vanishes on $f(\Lambda)$ has zero derivative on
$B$, then
\[
       J\subset\mathfrak q
       \quad\hbox{for every prime }
       \mathfrak q\subset\ker x.
       \tag{QX.5}
\]
\end{lemma}
\begin{proof}
Put $\mathfrak p=\ker x$ and
$\mathfrak a=\ker(xf)$.  These characteristic-zero points have
residue fields finite over $\operatorname{Frac}O$.  To check the
point without a hidden continuity assumption, present $B$ as a
quotient of $O[[X_1,\ldots,X_s]]$.  Each $x(X_i)$ is algebraic and
has absolute value less than one: otherwise a polynomial relation
with unit constant term would give a unit in the maximal ideal of
the power-series ring.  A distinguished polynomial annihilates
each $x(X_i)$, and division by these polynomials reduces the entire
quotient to a finite $O$-module.  The map is consequently continuous
and has image in a finite extension.  Squaring these polynomial
relations proves the corresponding finite coefficient statement
for a dual-number lift.  This is the point comparison of DX5.

The regular local ring $\Lambda_{\mathfrak a}$ has dimension $N$.
The tangent hypothesis makes restriction injective on derivations
at the point.  Since its residue field is separable over
$\operatorname{Frac}O$, the tangent space of
$\Lambda_{\mathfrak a}$ has dimension $N$.  The cotangent space
of $B_{\mathfrak p}$ therefore has dimension at most $N$.

On the other hand $(B/J)_{\mathfrak p}$ has dimension at least
$N$.  Indeed every minimal prime of the finite
$\Lambda$-algebra $B/J$ contracts to zero: minimal primes are
associated primes, and a nonzero element of $\Lambda$ is a
nonzerodivisor by torsion-freeness.  On a component through the
point we therefore have a finite integral extension of the domain
$\Lambda$.  Going down over the integrally closed ring $\Lambda$
lifts a full prime chain under $\mathfrak a$ to one ending at the
chosen point.  Its length is $N$.

Dimension is bounded by embedding dimension.  Hence
$B_{\mathfrak p}$ is a regular local ring of dimension $N$, and
its quotient by $J_{\mathfrak p}$ has the same dimension.  A regular
local ring is a domain; every nonzero ideal in it lowers dimension.
Thus $J_{\mathfrak p}=0$.  For $a\in J$ choose
$s\notin\mathfrak p$ with $sa=0$.  If
$\mathfrak q\subset\mathfrak p$ is prime, $s\notin\mathfrak q$
and therefore $a\in\mathfrak q$.  This proves (QX.5), including
the case $N=0$.
\end{proof}
% END QX2

% BEGIN QX3
\begin{proposition}[QX3: the actual characteristic-zero support leaf]
\label{claim:QX3}
Let $B,\Lambda$ be actual ordinary objects for a level $V$ and
residual datum $\bar r$ with \texttt{V0ok}.  Let $J$ have the
finite and torsion-free properties in QX2.  Suppose that an actual
point $x:B\to\overline{\mathbb Q}_\ell$ kills $J$ and is represented
by a continuous $\rho:G_L\to\operatorname{GL}_n
(\overline{\mathbb Q}_\ell)$ with the ordered flags
$\psi_{v,i}=x(\psi^B_{v,i})$.  Assume $n\ge2$,
\texttt{IsAutCsdA}$\,\iota\rho$, cyclotomic irreducibility,
unramifiedness at the auxiliary place, and integers $m_{v,i}$ and
$k>0$ such that
\[
 m_{v,i}\ge m_{v,j}+2\ (i<j),\qquad
       (\psi_{v,i}\chi^{-m_{v,i}})^k=1
                 \quad\hbox{on }I_v.
       \tag{QX.6}
\]
Then every prime below $\ker x$ contains $J$.
In particular this supplies the entire fixed-weight support
conclusion of the source's \texttt{eigenbasis\_rank\_four}, without
requiring its unused residual self-twist and large-order conditions.
\end{proposition}
\begin{proof}
We prove the tangent hypothesis of QX2.  Take a dual-number lift
$x'$ constant on $\Lambda$.  At the absolutely irreducible residual
point, its determinant law is the determinant of a matrix
representation $\rho'$ lifting $\rho$, unique up to conjugacy.
This is the split absolutely irreducible case of Chenevier,
\emph{The $p$-adic analytic space of pseudocharacters of a profinite
group and pseudorepresentations over arbitrary rings}, Theorem 2.22,
\url{https://gaetan.chenevier.perso.math.cnrs.fr/articles/determinants.pdf}:
first pass to the Cayley--Hamilton quotient, then lift matrix units
over the henselian ring $E[\epsilon]/\epsilon^2$.  The resulting
algebra is $M_n(E[\epsilon]/\epsilon^2)$.  Continuity is preserved:
choose finitely many group matrices forming a basis modulo
$\epsilon$; the trace pairing expresses the coordinates of every
$\rho'(g)$ through the continuous determinant traces at finitely
many products of group elements.  QX2 permits a single finite $E$
for this argument.

Write $\rho'(g)=(1+\epsilon y(g))\rho(g)$.  Burnside's theorem
makes the group matrices span $M_n(E)$ at the given irreducible
point.  Nakayama then makes their lifts span
$M_n(E[\epsilon]/\epsilon^2)$; finitely many already suffice.
The trace pairing on this matrix algebra is nondegenerate, as
the matrix units show.  Thus the actual universal equations
$\operatorname{tr}(\rho'(h)Q)=0$ for all $h$ force $Q=0$,
both for an ordinary product and for each prescribed inertial
matrix $\rho'(g)-1$.  These arguments derive the strong matrix
identities at this point from the weak \texttt{Pair} relations.
The ordinary identities give the ordered flags of $\rho'$
with characters $x'(\psi^B_{v,i})$.  For completeness,
on inertia kill the finitely many characters in (QX.6), and choose
an element whose cyclotomic value has infinite order.  Its ordered
eigenvalues are distinct.  Lift their spectral projectors to the
dual numbers.  In the resulting basis, the ordinary product
identities, with this element in all but one factor, force each
entry below the prescribed flag to vanish.  This gives an actual
flag over the dual numbers, not just a factorization of its
characteristic polynomial.

The \texttt{P} identities equate the characteristic polynomials,
and hence the characters, of $\rho'$ and its polarized conjugate
dual.  In the same matrix algebra the trace pairing makes equality
of their characters an intertwining isomorphism.  Lifting its
residual matrix gives an invertible form $P'$ for $\rho'$; this is
also the uniqueness assertion in the split Cayley--Hamilton
construction.  The
first derivative of this identity says that the class of $y$ is
fixed by the polarized complex-conjugation action.  Since $2$ is
invertible in $E$, inflation--restriction for
$\Gamma/\Gamma_0$ identifies these fixed classes with
$H^1(\Gamma,R)$, where $R$ has the action (QX.3).  Thus there is a
continuous cocycle $z$ extending the class of $y$.  Changing $y$
by a coboundary, if necessary, makes this an equality on
$\Gamma_0$.

We verify every local finite condition.  At places over $\ell$,
constancy on $\Lambda$ fixes all the inertial diagonal characters.
After a finite local extension the finite-order factors in
(QX.6) disappear.  The successive strictly upper-triangular
characters are $\chi^{m_{v,i}-m_{v,j}}$ times unramified
characters, with $m_{v,i}-m_{v,j}\ge2$.  Successive extensions of
the ordered characters make the base flag crystalline: at each
step its extension module has all these positive differences,
so local duality and the Bloch--Kato exponential give
$H^2=0$ and $H^1=H^1_f$.  This is an induction on the length of
the flag, beginning with the crystalline rank-one characters.

Here is the dimension argument for the tangent, which avoids
assuming surjectivity for a finite local condition.  Let
$\mathfrak n$ be the strictly upper-triangular adjoint module,
and $\mathfrak b$ the upper-triangular one.  They are crystalline
subquotients of the adjoint of this crystalline flag, and
\[
 0\longrightarrow\mathfrak n\longrightarrow\mathfrak b
       \longrightarrow E^n\longrightarrow0.
 \tag{QX.7}
\]
Every graded character of $\mathfrak n$ has exponent at least
two.  Consequently $H^0(\mathfrak n)=H^2(\mathfrak n)=0$ and
$H^1(\mathfrak n)=H^1_f(\mathfrak n)$, of dimension
$\dim_E t_{\mathfrak n}$, where
$t_M=D_{\rm dR}(M)/\operatorname{Fil}^0$.
The long exact sequence of (QX.7) shows that the inverse image
of $H^1_{\rm unr}(E^n)$ in $H^1(\mathfrak b)$ has dimension
\[
 \dim_E t_{\mathfrak n}
       -(n-h^0(\mathfrak b))+n
       =\dim_E t_{\mathfrak n}+h^0(\mathfrak b).
 \tag{QX.8}
\]
Functoriality puts $H^1_f(\mathfrak b)$ inside that inverse
image.  The Bloch--Kato dimension formula gives its dimension
as $\dim_E t_{\mathfrak b}+h^0(\mathfrak b)$, and
$t_{\mathfrak b}=t_{\mathfrak n}$ since the quotient has weight
zero.  Thus the two spaces coincide.  Constancy on $\Lambda$
puts the tangent class in this inverse image.  Its image in
the full adjoint is therefore finite.  This proves the local
claim with exactly the gaps in (QX.6).  Restriction followed by
corestriction descends it from the finite local extension.

At a non-$\ell$ place in the prescribed level, the universal
inertial relations make $\rho'$ trivial on the stated finite-index
subgroup $I_p\cap\ker\bar r$.  The cocycle is therefore zero there;
averaging over the remaining finite quotient makes its class
unramified.  Outside the level this is built into $\Gamma$.
At the auxiliary place $v_0$, the base representation is
unramified.  Choose a stable lattice over a sufficiently large
finite coefficient field.  Its Frobenius characteristic
polynomial reduces, by the determinant's residual identity,
to $(T-a)^n$ for one nonzero $a$; the residual matrix need not
itself be scalar.  All the residual eigenvalue ratios of its
adjoint are nevertheless $1$.  Moreover
$q_{v_0}\not\equiv1\pmod\ell$ by \texttt{V0ok}.  Hence
$\operatorname{Ad}(\rho(\operatorname{Frob}_{v_0}))-q_{v_0}$
is invertible: its determinant reduces to
$(1-q_{v_0})^{n^2}$.  The tame relation forces the infinitesimal
inertia value to be zero; the wild pro-$p$ subgroup has no
continuous additive image in an $\ell$-adic vector space for
$p\ne\ell$.  This proves the required unramified condition at
$v_0$ as well.  KB6 and BX transport all these checks to the
original valuations and \texttt{Af.IsBK} cochain convention.

QX1 now makes $z$, and hence $y$, a coboundary.  Thus $\rho'$ is
conjugate to the constant lift of $\rho$.  Its ordered diagonal
characters are constant too: the distinct inertia projectors
fix their order, so conjugacy cannot permute them.  The determinant
and every ordered local character of $x'$ agree with the constant
ones.  We do not apply the source's universal property to a
dual-number field, which is not an object of $\mathcal C_O$.
Instead KX7--KX8, or Chenevier's Remark 3.5 followed by the
ordinary relation quotient, shows that $B$ is topologically
generated by the determinant coefficients and the ordered
local character values.  Both point maps are continuous by
the finite coefficient argument in QX2.  Equality on these
generators makes $x'$ itself constant.
QX2 proves the result.  Notice that the proof required membership
$x(J)=0$; it does not create a modular quotient or an automorphic
point on one.
\end{proof}
% END QX3

\subsection{Two actual links, with one coefficient threshold}

Fix an actual $\texttt{Fit}+\texttt{Deep}$ tower $T$ for the
original $D$.  Thus $F/\mathbb Q$ is totally real and abelian,
$K/\mathbb Q$ is imaginary quadratic, $L=FK$, and $L/F$ is
unramified at finite places.  The primes in the original $S$,
including $2$ and $\ell$, split in $K$.  The residual characters
and the chosen $\ell$-adic decomposition groups are the ones
in $T$.  In particular \texttt{Deep} says that
$\psi|_{G_F}$ has finite order prime to $\ell$.
Write $\theta_i=\chi_i(D,T.\texttt{pr})$, so their labels include
the source's case-dependent permutation.

% BEGIN QX4
\begin{proposition}[QX4: construction of the two literal links]
\label{claim:QX4}
There is a coefficient ring $O_1$ such that, for every permitted
$O'\supset O,O_1$, there are two literal objects
$\mathcal L,\mathcal L':\texttt{Link}(T,O')$ with identical
characters $\kappa,\beta,t$ and with different auxiliary
rational primes.  They use the original tower, its original
places and its original residual-character labels.  The choice
of $O_1$ is independent of any later \texttt{Space}.
\end{proposition}
\begin{proof}
We give the character construction first, before choosing either
auxiliary prime.  Write $\omega$ for the Teichmuller part of
$\chi$.  The \texttt{Ctx2}/\texttt{Ctx3} and \texttt{Good}
identities make
\[
       \eta=(\omega^{e+3}\xi\psi_{\ell'})^{-1}
       \tag{QX.9}
\]
an even finite character of $G_{\mathbb Q}$ of order prime to
$\ell$, unramified outside $S$.  Here $\psi_{\ell'}$ is the
prime-to-$\ell$ part of $\psi$, obtained by an integer exponent
which is zero on its $\ell$-primary part and one on its
prime-to-$\ell$ part.  Deepness gives
$\psi_{\ell'}=\psi$ on $G_L$.

We need a finite character $\kappa_0$ of $G_K$ with
$\kappa_0\kappa_0^c=\eta|_{G_K}$ and the same permitted
ramification.  This has a direct ray-class construction.
Represent $\eta$ by an even Dirichlet character modulo $M$
supported on its ramified primes.  Choose, over each such prime,
one of the two primes of $K$; the splitting and Chinese
remainder theorem give $\mathcal O_K\to\mathbb Z/M\mathbb Z$.
For a principal ideal $(a)$ prime to $M$, evaluate the
Dirichlet character on this image of $a$.  The value is
independent of its generator: $2$ splits in $K$, so the units
of $\mathcal O_K$ are $\{\pm1\}$, and $\eta$ is even.
This character on principal ideals extends to the group of
ideals prime to $M$ with values in the divisible group of
prime-to-$\ell$ roots of unity.  It is trivial on principal
ideals congruent to $1$ modulo $M$, so it factors through the
finite ray class group.  Class field theory gives $\kappa_0$.
For every ideal $\mathfrak a$, the equality
$\mathfrak a\overline{\mathfrak a}=(N\mathfrak a)$ proves the
norm identity on Frobenius elements, hence everywhere by
Chebotarev.  Inverting the character if the opposite Artin
normalization is used gives the stated identity in the source's
Frobenius convention.

The unique square root in $1+\ell\mathbb Z_\ell$ gives the
continuous character
$\gamma=(\chi/\omega)^{1/2}$.  It is trivial on $\ker\chi$.
Next use the chosen prime of $K$ over $\ell$.  If $h_K$ is the
class number, a principal generator of $\mathfrak a^{h_K}$
for ideals away from $\ell$ gives, under the corresponding
embedding $K\hookrightarrow\mathbb Q_\ell$, the ray-class
character with infinity exponent $h_K(\ell-1)$ at that
embedding and zero at its conjugate.  Passing to the inverse
limit of the $\ell$-power ray class groups gives a continuous
$\mu:G_K\to1+\ell\mathbb Z_\ell$ and an integer $d>0$ with
\[
 \mu\mu^c=\chi^d,\qquad
 \mu|_{I_{\mathfrak l}}=\chi^d,\qquad
 \mu|_{I_{\bar{\mathfrak l}}}=1 .
 \tag{QX.10}
\]
One can take a positive multiple of $h_K(\ell-1)$ for $d$.
The ideal description proves these identities on principal
ideles; finite class field theory gives them on the inverse
limit.  This is a construction of the character, not an
assumption of a $\mathbb Z_\ell$-extension with additional
local behavior.

Set $\pi_0=\mu^2\chi^{-d}$.  Then $\pi_0^c=\pi_0^{-1}$, its
image lies in $1+\ell\mathbb Z_\ell$, and it is unramified
outside $\ell$.  There is a continuous additive
$t:G_K\to\mathbb Z_\ell$ with $\pi_0=(1+\ell)^t$.
Choose $a\in\overline{\mathbb Q}_\ell$ with
$a^{2d}=1+\ell$ and $|a-1|_\ell<1$, and put $\beta=a^t$.
The continuous $\mathbb Z_\ell$-power is defined as the limit
of integral powers; $a^{\ell^r}\to1$ proves that it exists.
It has image in a single finite coefficient field and satisfies
\[
 \beta^{2d}=\pi_0,\quad \beta^c=\beta^{-1},\quad
 \ker t\subset\ker\beta,\quad \bar\beta=1,\quad
 (\beta^2\chi^{-b_v})^d=1\ \hbox{on }I_v,\quad b_v\in\{1,-1\}.
 \tag{QX.11}
\]
In particular the exponents $b_v$ are odd.
Restrict these characters to $G_L$ and define
\[
        \kappa=\kappa_0\,\gamma^{-(e+4)}
                              \beta^{\,e\bmod2}.
        \tag{QX.12}
\]
Using (QX.9), $\gamma^2\omega=\chi$, and Deepness, direct
multiplication gives
\[
 \kappa(g)\kappa(cgc^{-1})\,\vartheta(g)\chi(g)^3=1,
 \quad \beta(g)\beta(cgc^{-1})=1,\quad
                         t(cgc^{-1})=-t(g).
 \tag{QX.13}
\]
At each place over $\ell$, (QX.11) shows that a positive power
of $\kappa\chi^{-a_v}$ is trivial for an \emph{integer} $a_v$:
the numerator $-(e+4)+b_v(e\bmod2)$ is even.
Away from $\ell$, $t$ and $\beta$ are unramified and
$\kappa$ has bounded prime-to-$\ell$ inertial order.
All three characters are trivial on the required
$\ell$-primary torsion of the local abelianization.  For
$\kappa_0$ this follows from its prime-to-$\ell$ order;
for $\gamma$ and $\beta$ it follows from their torsion-free
one-parameter images.  This also proves the exact
\texttt{mu} clause for the chosen decomposition groups.

Choose $O_1$ containing the finite character values, $a$, and
all embeddings of $\mathcal O_L$ into
$\overline{\mathbb Q}_\ell$.  The latter is a finite requirement
because $\mathcal O_L$ is finitely generated over $\mathbb Z$.
All character images then lie in $O_1^\times$.  Their
continuity is adic continuity over every $O'\supset O_1$.
This verifies the \texttt{large} field before choosing any
space or auxiliary prime.

We construct the auxiliary place.  The \texttt{g2} condition
in $\texttt{Fit}$ says that the residual cyclotomic character
is not in the subgroup generated by the two torus-root
characters.  Every ratio $\bar\theta_i/\bar\theta_j$ belongs
to that subgroup, using
$\theta_0\theta_3=\theta_1\theta_2$.
Duality for the finite abelian image of these characters gives
an element $g_0\in G_L$ with all $\bar\theta_i(g_0)$ equal
and $\bar\chi(g_0)\ne1$.  Indeed, failure of this conclusion
would put $\bar\chi$ in the annihilator of the common kernel,
which is exactly the subgroup generated by those ratios.

Choose an open normal subgroup on which these residual
characters are trivial.  Chebotarev, applied to its finite
Galois quotient and the class of $g_0$, gives a place $v_0$
of $F$ of residue degree one over a rational prime $p$,
split in $L$, with this Frobenius class.  The prime can avoid
any preassigned finite set.  Take that set to include the
original level, the discriminants, $2,\ell$ and the
ramification of the finite characters.  Then the residual
characters coincide on the whole decomposition group at
$v_0$, inertia is killed there, and $p\not\equiv1\pmod\ell$.
The last condition also proves $\mu_\ell\not\subset L$.

Take $S_L$ to be the places of $F$ over the original $S$
together with $v_0$, with the source's specified
$\ell$-adic places.  All are split in $L$; $[F:\mathbb Q]$
is even by Fit.  These are exactly the arithmetic conditions
of the \texttt{HidaSetting}.  KX2 supplies its level with
$V.S'=S_L$.  The characters
\[
 \bar r_i=\overline{\kappa\,\theta_i|_{G_L}},
 \qquad \bar\psi_{v,i}=\bar r_i|_{G_v},\qquad
                        \texttt{th}_{v,i}=1
 \tag{QX.14}
\]
descend to that quotient.  They have open kernels and remain
pairwise distinct by Fit; (QX.13) gives their determinant
polarization.  The previous paragraph is exactly
$\bar r.\texttt{V0ok}$.  All fields of \texttt{Link} are now
accounted for.

Repeat only the Chebotarev choice, adding $p$ to the excluded
set.  This gives $p'\ne p$ and hence two links with the
\emph{same} $\kappa,\beta,t$ and the same $O_1$.
No embedding, original residual label, or original
$\ell$-adic place has changed.
\end{proof}
% END QX4

% BEGIN QX5
\begin{proposition}[QX5: the actual determinant comparison maps]
\label{claim:QX5}
For either link $\mathcal L$ in QX4, every actual
$X:\texttt{Space}(T.\texttt{fr}\,\varnothing,O\to O')$ and every
integer $a\ge0$, there is an actual ordinary pair over $X.B$.
It supplies a universal ordinary determinant ring $B$ by KX9.
For \emph{every} such universal ring there is a unique comparison
map $\phi_a:B\to X.B$ satisfying the literal
$\mathcal L.\texttt{IsPhi}\,X\,a\,B\,\phi_a$:
\[
 \begin{split}
 \phi_a\bigl(\det(T-D_B(g))\bigr)
  &=\operatorname{scaleRoots}_{\kappa(g)\beta(g)^a}
                    \bigl(X.\operatorname{cp}(g|_F)\bigr),\\
 \phi_a(\psi^B_{v,i}(g))
  &=\kappa(g)\beta(g)^a\,X.\psi_{v,i}(g|_F).
 \end{split}
 \tag{QX.15}
\]
The first line denotes the characteristic polynomial of the
universal determinant at $g$, rather than a matrix realization
of $D_B$.  Both maps land in the original $X.B$, with its
specified injection into $X.A$.
\end{proposition}
\begin{proof}
In $X.A$ take the actual matrices
$\rho_a(g)=\kappa(g)\beta(g)^a X.P.r(g|_F)$.
Their characteristic coefficients lie in $X.B$ by the first
formula in (QX.15).  They are adically continuous and reduce
to (QX.14).  We first explain why this gives a determinant
\emph{law} over $X.B$, even in residue characteristic three.

For finitely many matrices $A_1,\ldots,A_s$, the integral
Amitsur identity expresses every coefficient of
$\det(1-\sum Z_iA_i)$ as an integer polynomial in the
characteristic coefficients of products of the $A_i$.
One form of the identity is the formal product of
$\det(1-Z_wA_w)$ over primitive cyclic words, in Lyndon
order; only finitely many words enter any fixed degree.
It follows by the unique nonincreasing Lyndon factorization
of a word and multiplicativity of determinant.  In the
present application each product $A_w$ is another group
matrix.  Consequently all these mixed coefficients belong
to $X.B$.  They define the determinant law there, and the
injection $X.B\hookrightarrow X.A$ verifies its homogeneity,
multiplicativity and all base-change identities.  This is
also the integral determinant descent in Chenevier,
Lemma 1.12 and Corollary 1.14.  It uses no division by $4!$.

Descent from $G_L$ to the level quotient needs a small
topological precaution.  The ring $X.A$ is only finite type
over $X.B$; its adic topology need not be separated, so we
do not infer that the kernel of the matrix representation
is closed.  Instead take the subgroup consisting of those
$u$ for which all characteristic coefficients at $gu$
equal those at $g$, for every $g$, and the corresponding
mixed coefficients do also.  The just-proved word identities
make this the normal kernel of the determinant law.
It is closed because its coefficients take values in the
separated ring $X.B$ and are adically continuous.  It
contains every inertia subgroup killed by $X.P.r$ and
$\kappa\beta^a$, so it contains their closed normal closure,
which is the kernel of $G_L\twoheadrightarrow\Gamma_0$
in KX2.  The determinant therefore descends to $\Gamma_0$.
Its continuity descends through this quotient map.

For the ordered local characters, the upper-triangular
matrices given by $X.P.n_v$ show directly that
$X.\psi_{v,i}(gh)=X.\psi_{v,i}(g)X.\psi_{v,i}(h)$ in $X.A$,
hence in $X.B$ by injectivity.  Their residual values are
nonzero, so they are units of the local ring $X.B$.
Multiplying by $\kappa\beta^a$ gives the second line of
(QX.15), including continuity, residual values and the
prescribed local torsion.

We check the remaining weak \texttt{Pair} relations.
The symplectic multiplier of $X.P.r$ and (QX.13) identify
the characteristic polynomial at a group element with
its polarized conjugate dual; $\beta\beta^c=1$ makes this
independent of $a$.  This is exactly \texttt{P}.
The actual ordered upper-triangular matrices make the
ordinary product zero, so multiplying by any group matrix
and taking its trace gives \texttt{Ord}.

For \texttt{I}, let $g$ be in one of the specified non-$\ell$
inertias and in all residual-character kernels (QX.14).
Choose an integer $m$ prime to $\ell$ killing
$\kappa$ on this inertia.  Here $\beta=1$.  The element
$g^m$ has original residual matrix $1$, so the defining
kernel condition of $T.\texttt{fr}\,\varnothing$ gives
$\rho_a(g)^m=1$.  Its characteristic polynomial, with
coefficients in $X.B$, reduces to $(T-1)^4$.
This already forces $\rho_a(g)=1$, including over $X.A$.
Indeed write $T^m-1=(T-1)Q(T)$, with $Q(1)=m$ a unit.
The resultant of $Q$ and that characteristic polynomial is
a unit of $X.B$, by reduction modulo its maximal ideal.
Bezout and Cayley--Hamilton make $Q(\rho_a(g))$ invertible;
multiplication by $(\rho_a(g)-1)$ then gives the result.
Its trace version is the required \texttt{I} relation.

We have constructed an actual \texttt{Pair} over the ring
$X.B\in\mathcal C_{O'}$.  Thus the nonempty-pair antecedent
in KX9 is genuinely satisfied, rather than discarded.
The universal property gives $\phi_a$, and evaluating
the determinant and characters gives precisely (QX.15).
Conversely those formulas determine their full polynomial
law by the word identity and determine every ordered
character.  Universality in the actual category
$\mathcal C_{O'}$ proves uniqueness.  Every construction
works for all $a\ge0$ and both links on the same space.
\end{proof}
% END QX5

\subsection{The ordinary ideal properties that are actually used}

For each link we now suppose an ideal $J_0\subset B$ has the
literal rank-four \texttt{HidaPropsAlt} fields.  We list the
fields to keep the boundary explicit:
\begin{enumerate}
\item \texttt{ring}: $B/J_0$ is finite and torsion-free over
$\Lambda$, with the specified polarization, ordinary,
inertial-type and auxiliary-unramified identities.
\item \texttt{dense}: for every strictly increasing integral
weight, the source's good arithmetic points killing $J_0$
separate $B/J_0$.
\item \texttt{pt}: an actual continuous attached regular
algebraic conjugate-self-dual cuspidal representation with
\texttt{WitA}, the stated auxiliary unramifiedness and
ordered local generic flags belongs to this quotient.
\item \texttt{end1}: every arithmetic point on the quotient
has the exact semisimple decomposition into consecutive
cyclotomic strings, and the stated strict Frobenius-root
bound when those strings all have length one.
\item \texttt{end2}: every irreducible arithmetic point is
actually conjugate-self-dual cuspidal and has
\texttt{WitOff} away from the auxiliary place.
\item \texttt{twist}: the good arithmetic points are stable
under the specified finite $\ell$-power-order polarized
twists, trivial on the prescribed local kernels.
\end{enumerate}
These are assumptions on an actual ideal in an actual
universal ring; a citation to ordinary automorphic forms
without their exact local types does not supply them.
The separate characteristic-$\ell$ support assertion
\texttt{P1core} will be stated at QX8.  No \texttt{P3core}
field is used in this chapter.

% BEGIN QX6
\begin{proposition}[QX6: integral twist invariance of the image ideal]
\label{claim:QX6}
Assume the \texttt{dense} and \texttt{twist} fields above.
For the maps of QX5,
\[
              \phi_a(J_0)X.B=\phi_0(J_0)X.B
                    \qquad(a\ge0).
              \tag{QX.16}
\]
Here equality is equality of ideals, before evaluating any
point and over the original coefficient ring $O'$.
\end{proposition}
\begin{proof}
The character $\beta$ descends to $\Gamma_0$, since it kills
the closed inertia kernel of the level quotient.  It is
conjugate-inverse, is $1$ residually, and is trivial on the
specified non-$\ell$ inertia and the local torsion.
Twisting the universal determinant and every ordered
character by such a character $\gamma$ therefore produces
another actual \texttt{Pair}.  Universality gives a map
$\tau_\gamma:B\to B$ for coefficient-valued $\gamma$,
with inverse $\tau_{\gamma^{-1}}$.  The coefficient of
degree $4-j$ at $g$ is multiplied by $\gamma(g)^j$.

We prove that $\tau_\gamma(J_0)\subset J_0$ without any
incorrect assertion that finite-order characters converge
pointwise to $\gamma$.  Fix $r\ge1$.  The image of $\gamma$
modulo $\mathfrak m_B^r$ is a finite abelian $\ell$-group
$Q$: it lies in $1+\mathfrak m_B$, whose successive
quotients are finite additive groups of characteristic
$\ell$.  Its quotient map $q:\Gamma_0\to Q$ has open
kernel and has all the same polarization and local-kernel
identities.  The group algebra $B[Q]$ is a finite free,
complete local $O'$-algebra with residue field equal to
that of $B$.  Its maximal ideal is generated by
$\mathfrak m_B$ and $[u]-1$, since a finite $\ell$-group
has local group algebra in characteristic $\ell$.
Twist by the tautological character $g\mapsto[q(g)]$.
This gives an actual map $\tau_Q:B\to B[Q]$.

Take $b\in J_0$ and write $\tau_Q(b)=\sum_{u\in Q}b_u[u]$.
For any good arithmetic point $x$ killing $J_0$ and any
character $\lambda:Q\to\overline{\mathbb Q}_\ell^\times$,
the \texttt{twist} field says
\[
                  \sum_{u\in Q}x(b_u)\lambda(u)=0.
                  \tag{QX.17}
\]
These twists have finite $\ell$-power order, preserve the
ordered-character ratios, and satisfy every stipulated
local condition.  Character orthogonality over the
characteristic-zero field makes each $x(b_u)$ zero.
The \texttt{dense} field then gives $b_u\in J_0$ for every
$u$.  Evaluating $[q(g)]$ at $\gamma(g)$ modulo
$\mathfrak m_B^r$ and using uniqueness of the twisted pair
shows
$\tau_\gamma(b)\in J_0+\mathfrak m_B^r$.
Krull separation of the complete Noetherian quotient
$B/J_0$ gives $\tau_\gamma(b)\in J_0$.
The inverse twist gives equality.

Apply this with the coefficient-valued character $\beta^a$.
Formula (QX.15) and QX5 uniqueness give
$\phi_a=\phi_0\tau_{\beta^a}$.  Equality of image ideals
follows.  This argument also proves the universal
$B$-valued twist-stability statement used in the source;
the finite group algebra, rather than a quotient requiring
division by $\ell$, is what makes it integral.
\end{proof}
% END QX6

% BEGIN QX7
\begin{proposition}[QX7: full local support at an actual seed]
\label{claim:QX7}
Under the stated rank-four \texttt{HidaPropsAlt} hypotheses,
put $J_{\mathcal L}=\phi_0(J_0)X.B$.  For every actual
$O'$-map $y:X.A\to\overline{\mathbb Q}_\ell$ for which
$T.\texttt{IsSeedAlt}\,\iota(X.P_y)$ holds, one has
\[
                  X.\texttt{LocZero}(J_{\mathcal L},y).
                  \tag{QX.18}
\]
This means that every prime of $X.B$ contained in
$\ker(y|_{X.B})$ contains the entire ideal
$J_{\mathcal L}$.
\end{proposition}
\begin{proof}
Unpack the seed as $(w,\kappa_s,\pi,a_s,k_s)$.
Its consecutive $w_i$ differ by even integers at least
four, $w_0+a_s$ is even, and
$(\kappa_s^2\chi^{-a_s})^{k_s}=1$ on every $\ell$-inertia.
Its automorphic representation is an actual regular
algebraic rank-four cusp form with \texttt{AllWin}, attached
to $\kappa_s X.P_y.r|_{G_L}$.  The full-matrix-span condition
holds already on $G_L\cap\ker\chi$.

Choose $a\ge0$ with $w_0+a$ even and use $\phi_a$ from QX5.
At each $v\mid\ell$, (QX.11)--(QX.12) give an exponent
$e_v\equiv a\pmod2$ for
$(\kappa\beta^a)^2$.  All the integers
\[
               m_{v,i}=(w_i+e_v)/2
               \tag{QX.19}
\]
therefore exist, and their differences are
$(w_i-w_j)/2$.  A common positive power makes
$(x_a(\psi^B_{v,i})\chi^{-m_{v,i}})$ trivial on inertia,
where $x_a=y|_{X.B}\phi_a$.  To obtain a single power on
all inertia, take the index of the normal core of the
open subgroup in $T.\texttt{Wts}$, and multiply by the
finitely many local exponents.  Thus (QX.6) holds with
all its gaps at least two.

The ratio $\lambda=\kappa\beta^a\kappa_s^{-1}$ is a
continuous algebraic Galois character.  In fact
$e_v-a_s$ is even, and a positive power of
$\lambda\chi^{-(e_v-a_s)/2}$ is trivial on every local
inertia.  Local class field theory makes it locally
algebraic there; global class field theory gives the
corresponding algebraic Hecke character.  This is the
rank-one algebraic-character theorem, with the integer
exponents just checked.  Twisting $\pi$ by that Hecke
character attaches
\[
                 \rho_a=(\kappa\beta^a)X.P_y.r|_{G_L}.
                 \tag{QX.20}
\]
The polarization (QX.13) and strong multiplicity one
identify the twist's conjugate with its contragredient,
so the result is conjugate-self-dual.  The same comparison
uses geometric Frobenius on both sides; scalar twisting
does not introduce an additional power of $\chi$.

We also retain the \texttt{WitA} part.  An algebraic
Hecke-character twist translates the center of every
strict window by its real norm exponent, preserving
\texttt{AllWin}.  Every other regular algebraic cusp
form attached to (QX.20) is the same form by strong
multiplicity one.  Finally (QX.19) implies
\texttt{PotUnl} at every place over $\ell$: an equality
between a positive power of two ordered characters
with one cyclotomic twist would, on an open inertia
subgroup, force $m_{v,i}-m_{v,j}=1$, which is impossible.
Thus both clauses of \texttt{WitA} follow, without a
general assertion that arbitrary attached forms have it.

Continuity here follows from the actual point, not from an
unstated clause of the seed.  Absolute irreducibility and
Burnside give finitely many matrices $\rho_a(g_j)$ spanning
the matrix algebra.  The nondegenerate trace pairing recovers
every matrix entry of $\rho_a(g)$ from the functions
$\operatorname{tr}(\rho_a(gg_j))$.  These are continuous
by the characteristic-coefficient continuity of $X$ and
the continuous small-point comparison of QX2; scalar
twisting is continuous as well.  Thus $\rho_a$ is continuous.
Its ordered flags are the actual
flags of $X.P_y$, and its cyclotomic irreducibility is
unchanged by the scalar twist.  At the auxiliary prime
the residual characters coincide, and the same
prime-to-$\ell$ finite-order matrix argument as in QX5
gives the required unramifiedness.  The
\texttt{PointTheoremBaseObject} field now says
$x_a(J_0)=0$: the ordinary weights are arithmetic by
local reciprocity, and their gaps exclude every forbidden
ratio $\psi_i=\psi_j\chi$.

Apply QX3 to $x_a$.  Given a prime
$\mathfrak q\subset\ker(y|_{X.B})$, its inverse image
under $\phi_a$ is a prime below $\ker x_a$.
It contains $J_0$, so $\mathfrak q$ contains
$\phi_a(J_0)X.B$.  QX6 identifies this ideal with
$J_{\mathcal L}$ and proves (QX.18).
In particular an actual seed supplies a nonzero point
on the appropriate support; it is not assumed at the
start of the fixed-weight argument.
\end{proof}
% END QX7

% BEGIN QX8
\begin{proposition}[QX8: the complete characteristic-$\ell$ support transport]
\label{claim:QX8}
Assume in addition the literal $B.\texttt{P1core}\,J_0$:
if an actual finite-field power-series point is
\texttt{NiceI}, matches the universal pair, and kills
$J_0$, every prime below its kernel contains $J_0$.
Then $J_{\mathcal L}$ satisfies the entire
\texttt{p1} clause of \texttt{PotentialModularityTower}
on the original $X$.  No further residual-image or
generic-point hypothesis is required.
\end{proposition}
\begin{proof}
Take the exact test in that clause: a finite field $k$,
an $O'$-map $y:X.A\to k[[t]]$ whose constant term kills
$\mathfrak m_{X.B}$, with $y(J_{\mathcal L})=0$,
the prescribed local character separations on every
$T.U_1(v)$, and full matrix span over $k((t))$.
The constant-term condition makes $O'\to k$ local,
so the given residual field embeds in $k$.
Twisting by $\kappa$ and descending as in QX5 gives
an actual representation over $k[[t]]$ of $\Gamma_0$,
its symmetric polarization and its ordered characters.
Here the matrix descent itself is valid, since
$k[[t]]$ is separated.  Its scalar-extended determinant,
residual characters, inertia and flags give the actual
source \texttt{Point}; its map from $B$ is
$f=y|_{X.B}\phi_0$.

We verify all four \texttt{NiceI} fields rather than
assuming that full span already means \texttt{NiceI}.
First, restriction of the original irreducible
representation from $G_F$ to $G_L$ stays absolutely
irreducible.  Otherwise Clifford theory for this
index-two subgroup would make its trace zero on the
other coset.  Its reduction is the sum of the four
original residual characters, whose restrictions to
$G_L$ are distinct by (QX.14).  On any other coset,
linear independence of those four restricted
characters forbids an identically zero sum with their
nonzero coset values as coefficients.  This contradiction
works in every odd characteristic, including three.

The local separation condition gives, for each ordered
pair $i\ne j$, an element of $\Gamma'$ at which the
residual characters agree but their power-series
characters do not.  The ratio there lies in
$1+tk[[t]]$ and has infinite order; this group is
torsion-free, by the Frobenius map for the $\ell$-part
and the simple-root identity for the prime-to-$\ell$
part.  Hence the kernel of the ratio modulo its finite
constant torsion is closed and not open.
Avoiding finitely many such kernels in the compact
local group gives one element $u\in\Gamma'$ whose
four eigenvalue ratios all have infinite order.
Every positive power of $\rho(u)$ has four distinct
eigenvalues.  In particular local characters are
distinct and no ratio is $\chi$, proving \texttt{N1}.

Suppose restriction to $\Gamma'$ were reducible.
Clifford theory gives a sum of conjugate simple
constituents.  Their multiplicity is one, since a
positive power of $\rho(u)$ still has simple spectrum.
The number of constituents is therefore $2$ or $4$.
The quotient $\Gamma_0/\Gamma'$ is abelian; a
nontrivial quadratic character of the resulting
permutation quotient gives a quadratic self-twist.
Its kernel is open: stabilizers of subspaces are
closed by continuity, and they have finite index.
Because $\ell$ is odd, this quadratic character
of a closed subgroup of $\mathbb Z_\ell^\times$
is trivial on its pro-$\ell$ part, so factors through
the residual cyclotomic character.  Reduction of
the characteristic-polynomial identity permutes
the four residual characters.  This is excluded
by the exact \texttt{g3'} field of Fit, transported
through $\kappa$.  Thus \texttt{TypeI} holds.

We give the rank-four regular-coset argument needed
for \texttt{N3}.  Work over an algebraic closure of
$k((t))$, and let $H$ be the Zariski closure of
the image of $\Gamma'$.  No open finite-index
subgroup of $\Gamma'$ can have commuting image.
If one did, its Zariski closure would be
commutative and contain $H^\circ$: finitely
many translates cover $H$, so one of them
contains its identity component.  The group
$H^\circ$ is normal in the Zariski closure of
$\Gamma_0$, since $\Gamma'$ is normal.
Its inverse image in $\Gamma'$ is closed
of finite index and contains a positive power
of $u$.  The four distinct eigenspaces of
that power are therefore permuted transitively
by $\Gamma_0$, making the representation
monomial.  On the kernel of the
permutation action, all four reduced diagonal
characters coincide.  To see this, reduce their
integral eigenvalues and compare with the four
given residual characters; the latter are
conjugation-invariant, whereas transitivity
permutes the former.  Independence of characters
gives the assertion.  Every residual ratio thus
factors through a transitive subgroup of $S_4$.
Such a subgroup is $C_4$, $C_2^2$, a dihedral
group of order eight, $A_4$, or $S_4$; its
one-dimensional characters have order at most
four.  This contradicts Fit's \texttt{ord} field.

The inverse image of $H^\circ$ is closed and of finite index,
hence open.  Every open subgroup is Zariski dense
in the relevant connected component.  A power
of $u$ lies there and generates a diagonal torus
whose four weights are distinct.  Choose a
cocharacter $z\mapsto\operatorname{diag}(z^{a_i})$
with distinct integers $a_i$; this avoids finitely
many proper hyperplanes in its cocharacter lattice.

Restriction to $H^\circ$ is semisimple by Clifford's
argument: conjugates of one simple submodule span
the whole irreducible representation.  The simple
summands have equal dimension and multiplicity
one by the simple spectrum just noted.  Dimension
one would make $H^\circ$ commuting, already
excluded.  Thus there is one four-dimensional
block or two two-dimensional blocks.  On any
component $gH^\circ$, a generic matrix $M$ therefore
has either all four diagonal entries and all
four diagonal cofactors nonzero, or the block
form
\[
                  M=\begin{pmatrix}0&A\\ B&0\end{pmatrix}
                  \tag{QX.21}
\]
with every entry of $A,B$ nonzero and
$\det A\det B\ne0$.  This follows from Burnside
on each simple block; nonvanishing is an open
condition on the irreducible variety $H^\circ$.

In the first case the characteristic polynomial of
$M\operatorname{diag}(z^{a_i})$ is separable over
the Laurent-series field.  Order the $a_i$ strictly.
The extreme linear and cubic coefficients have
nonzero leading terms, so its Newton polygon has
two outer edges of length one; any remaining
edge has length at most two.  A length-two edge
has no intermediate leading coefficient: the
possible pair sums jump strictly past its midpoint.
Its residual binomial is separable because
$\operatorname{char}k\ne2$.  In the block-swap
case the polynomial is
$T^4-c(z)T^2+d(z)$.
The largest term of $c(z)$ is the unique product
using the largest exponent in each block and
has nonzero coefficient.  Its exponent is
strictly more than half the exponent of $d(z)$.
The two roots in $T^2$ have distinct orders and
each has two distinct square roots.  This proves
separability in that case too.

The discriminant is therefore a nonzero regular
function on a suitable translate of that torus.
Powers of $u$ are Zariski dense in it.  Choosing
the generic $M$ inside any prescribed open
coset and then multiplying by a suitable power
of $u$ produces a separable characteristic
polynomial in that coset.  This proves the
source's \texttt{N3} on every nonempty open
intersection with $\Gamma'$.

Finally \texttt{N4} is the auxiliary-prime
calculation already used in QX3.  The
Frobenius characteristic polynomial reduces
to $(T-a)^4$ with $a\ne0$, while the residual
cyclotomic value is not $1$.  Hence the
operator expressing a cyclotomic-twisted
adjoint invariant has invertible determinant
over $k[[t]]$.  Clearing denominators proves
that its kernel over $k((t))$ is zero.

We have proved the full \texttt{NiceI} hypothesis
for this actual point.  Apply \texttt{P1core}
to $f$.  For a prime
$\mathfrak q\subset\ker(y|_{X.B})$, its inverse
image under $\phi_0$ is prime below $\ker f$.
It contains $J_0$, so $\mathfrak q$ contains
$J_{\mathcal L}$.  This is exactly the
required \texttt{LocZero}, with all original
finite-field and local-separation quantifiers.
\end{proof}
% END QX8

\subsection{Purity and the original-field automorphic conclusion}

We use the following precise classical results in the final transport.
Arthur--Clozel, \emph{Simple algebras, base change, and the advanced
theory of the trace formula}, Chapter III, Theorems 4.2 and 5.1,
give cyclic prime-degree cuspidal descent and base change, with
strong local compatibility; Chapter I, \S6 gives their local
tempered compatibility.  The primary text is
\url{https://www.claymath.org/library/cw/arthur/pdf/30.pdf}.
Harris--Lan--Taylor--Thorne, \emph{On the rigid cohomology of
certain Shimura varieties}, Theorem A (Corollary 7.14),
\url{https://virtualmath1.stanford.edu/~rltaylor/rigcoh.pdf},
attaches a continuous semisimple Galois representation, with the
geometric normalization above, to every regular algebraic cusp
form over a CM field.  No polarization is required for this
existence assertion.  We use it only to identify a finite-character
ambiguity in cyclic descent.

For a regular algebraic conjugate-self-dual cusp form over a CM
field, Caraiani's \emph{Local-global compatibility and the action
of monodromy on nearby cycles}, Theorems 1.1--1.2, and
\emph{Monodromy and local-global compatibility for $l=p$},
Theorem 1.1, give full Frobenius-semisimplified local compatibility
and temperedness at every finite place, for every rank $n\ge2$:
\url{https://www.ma.imperial.ac.uk/~acaraian/papers/lgc1.pdf} and
\url{https://www.ma.imperial.ac.uk/~acaraian/papers/lgc2.pdf}.
There is no Shin-regularity condition in these statements.
The word ``cohomological'' in Caraiani's statements matches
the actual regular algebraic cusp used here: Clozel's
\emph{Motifs et formes automorphes}, Lemma 3.14, identifies
its archimedean representation with one having nonzero
relative Lie algebra cohomology for the corresponding
algebraic coefficient, as already recorded in the CX
classical-input registry.
For the polarized forms obtained after descent, we use
BLGGT, Theorem 2.1.1(2)--(3), and \emph{Local-global
compatibility for $l=p$, II}, Theorem A,
\url{https://virtualmath1.stanford.edu/~rltaylor/lg2.pdf}.
The latter identifies the semisimplified Weil representations
at $\ell$ for every regular algebraic polarized cusp form;
we do not invoke its stronger, Shin-regular conclusion.
The needed vanishing of monodromy is proved below.

% BEGIN QX9
\begin{proposition}[QX9: purity and the complete automorphic return]
\label{claim:QX9}
Let $J_{\mathcal L}=\phi_0(J_0)X.B$ for one of the actual
links of QX4--QX5, and assume the six stated
\texttt{HidaPropsAlt} fields for $J_0$.
For every $y:X.A\to\overline{\mathbb Q}_\ell$ killing
$J_{\mathcal L}$ and satisfying
$T.\texttt{Wts}(2D.\texttt{wts4})$, the entire
\texttt{pur} clause of \texttt{PotentialModularityTower}
holds.  If, as in its \texttt{aut} clause, the original
$r:G_{\mathbb Q}\to\operatorname{GL}_4
(\overline{\mathbb Q}_\ell)$ is continuous, has
multiplier $D.\texttt{th}$, restricts to $X.P_y$ over
$F$, and is irreducible on $G_L$, then the literal
\texttt{IsAut4Alt}$\,\iota r$ holds, with the same
$\iota$ and for every imaginary quadratic field in
its definition.
\end{proposition}
\begin{proof}
Put $x=y|_{X.B}\phi_0$.  The weight calculation
in QX7, now with the even weights $2D.\texttt{wts4}$,
makes $x$ an actual arithmetic point.  The local
integer weights are
$D.\texttt{wts4}(i)+e_v/2$; here the exponent $e_v$
of $\kappa^2$ is even, since $a=0$.
Local reciprocity identifies their restrictions to
units with the required algebraic weight.  This
establishes the source's \texttt{Arith} condition,
not merely a Hodge-weight multiset.

Apply \texttt{end1} to $x$.  It supplies a continuous
semisimple $\rho$ with the actual determinant at $x$
and a characteristic-polynomial decomposition
\[
 \rho^{\rm ss}\simeq
 \bigoplus_i\ \bigoplus_{j=0}^{m_i-1}
                  \rho'_i\chi^{-j},
 \qquad d_i=\dim\rho'_i\ge1,\quad m_i\ge1.
 \tag{QX.22}
\]
All the $m_i$ are one.  Indeed choose stable lattices
in the finitely many continuous constituents.  Their
semisimplified reductions together give the four
distinct residual characters of $\kappa X.P$.
If $m_i\ge2$, a character of the reduction of
$\rho'_i$ and its cyclotomic twist both occur.
Their ratio is $\bar\chi$ or its inverse.
The common factor $\bar\kappa$ cancels, contrary
to Fit's exact \texttt{g2} exclusion of
$\bar\chi$ from the subgroup generated by the
original residual root characters.  This argument
uses reductions of stable lattices, not an
unjustified specialization of characteristic-zero
irreducible constituents.

Consequently \texttt{end1} gives, outside a finite
set of places $w$ of $L$, the strict inequalities
\[
        |\iota\alpha|<\operatorname{N}w\,
                           |\iota\alpha'|
        \quad\hbox{for every pair of Frobenius roots.}
        \tag{QX.23}
\]
The semisimple $\rho$ is unramified outside the
finite level set.  To check this from its determinant,
the universal law is constant on each killed
inertial coset.  Hence
$\operatorname{tr}(\rho(g)(\rho(u)-1))=0$ for
every $g$ and every killed inertia element $u$.
The semisimple image algebra is a product of matrix
algebras, and its trace pairing is nondegenerate in
characteristic zero.  Thus $\rho(u)=1$.

Take $H=G_L\subset G_{\mathbb Q}$ and let $M\ne0$
be the product of $\ell$, the rational primes under
these exceptional places, and the finitely many
level and discriminant primes.  If
$\sigma\in H$ is an arithmetic Frobenius of degree
$f$ at a rational prime $p\nmid M$, its image at
$w\mid p$ is the $(f/f_w)$-th power of the local
arithmetic Frobenius, where
$\operatorname{N}w=p^{f_w}$ and $f_w\mid f$.
The source \texttt{AttC} convention uses geometric
Frobenius, but (QX.23) is preserved on replacing
all roots by their inverses and exchanging the
ordered pair.  Raising it to the positive power
$f/f_w$ gives the bound $p^f$.
The characteristic polynomial of $\rho(\sigma)$
is that of $\kappa(\sigma)r(\sigma)$ by QX5.
This scalar cancels from every ratio.  Thus
\[
       |\iota z|<p^f|\iota z'|
       \quad(z,z'\hbox{ roots of }r(\sigma),\ z'\ne0),
       \tag{QX.24}
\]
with exactly the open subgroup, nonzero $M$, original
valuation, positive degree and Frobenius quantifiers
of \texttt{PureW}.  Neither irreducibility nor
continuity of the given extension $r$ was used here.

For the automorphic assertion, assume its additional
hypotheses.  The scalar twist
$\rho_L=\kappa r|_{G_L}$ is irreducible and
continuous.  The \texttt{end2} field at $x$ gives
an actual regular algebraic conjugate-self-dual
cuspidal $\pi_L$ attached to $\rho_L$.
In fact this field supplies \texttt{AutOffN},
including a window assertion away from the auxiliary
prime.  We use its actual cusp and the matched
stronger local theorems above.  This is a deliberate
simplification: two auxiliary primes remain available
for the integral support construction, but the
automorphic return does not need an unmatched
\texttt{WinOf} or \texttt{HypothesesH'} package.

The character $\kappa$ is algebraic by (QX.11)--(QX.12)
and the rank-one reciprocity argument of QX7.
Twist $\pi_L$ by its inverse Hecke character.
The resulting regular algebraic cusp $\pi_L^0$
attaches to the \emph{original} $r|_{G_L}$.
The twist also transports all local Langlands
parameters; it does not change the rank and so
introduces no extra Tate twist.

We spell out cyclic descent with the actual Galois
representation fixed.  Suppose $E/E_0$ is a cyclic
prime-degree step in $L/K$, and an actual cusp
$\pi_E$ attaches to $r|_{G_E}$.  This restriction
is irreducible, since its further restriction to
$G_L$ is.  For $\tau\in\operatorname{Gal}(E/E_0)$,
the attached representation is unchanged up to
isomorphism; strong multiplicity one gives
$\pi_E^\tau\simeq\pi_E$.
Arthur--Clozel III.4.2(d) supplies a cuspidal descent
$\pi_0$.  Its archimedean infinitesimal character
is the same at each extending complex embedding,
so it is regular algebraic.
Use HLTT to attach the continuous semisimple
$r_{\pi_0,\iota}$ without first presupposing
polarization of this chosen descent.
Its restriction to $G_E$ equals $r|_{G_E}$ by
Chebotarev and base change.  Both extensions of
this irreducible representation differ by a
character of $\operatorname{Gal}(E/E_0)$:
after choosing an intertwiner on $G_E$, Schur's
lemma makes the ratio of the two matrices a scalar
on every coset; these scalars form the required
finite continuous character.  Twist $\pi_0$ by
its finite Hecke character to obtain precisely
$r|_{G_{E_0}}$.  This proves descent without
silently changing the original Galois representation.
Iterating reaches the original imaginary quadratic
$K$, since $L/K$ is abelian.

The multiplier also descends correctly.  Set
$\theta=D.\texttt{th}$.  The symplectic pairing
and complex conjugation give
$r^c\simeq r^\vee\theta$ over every imaginary
quadratic field.  In our normalization the
Galois representation of $\pi^\vee$ is
$r_\pi^\vee\chi^{-3}$.  Thus strong multiplicity
one makes the descended cusp essentially
conjugate-self-dual with the algebraic Hecke
character corresponding to $\theta\chi^3$.
This character comes from $\mathbb Q$.  If its
real-place sign is not the required rank-four
sign $+1$, multiply it by the quadratic Hecke
character $\eta_{K'}$ of the imaginary quadratic
field in question.  That character changes the
real sign and is trivial after norm pullback
to $K'$; its Galois avatar is trivial on
$G_{K'}$.  Hence the automorphic polarization
and its Galois multiplier on $G_{K'}$ are
unchanged.  We do not identify a Hecke infinity
sign with the value of its uncorrected
cyclotomic avatar at complex conjugation, nor
assert global avatar equality after this
correction.  This proves all
the hypotheses of the polarized Galois and local
compatibility theorems used below.  The multiplier
is not replaced by its restriction or by a
character with a merely similar infinity type.

Now let $K'$ be any imaginary quadratic field
such that $r|_{G_{K'}}$ is irreducible.  The
representation $r$ is unramified outside a finite
set: outside the original residual ramification
and the discriminant of $F$, inertia lies in
$G_F\cap\ker\bar r$, where the actual
$X.\texttt{pair.kill}$ makes it trivial.
Choose another imaginary quadratic $K_0$ ramified
at an odd prime $q$ outside this set and outside
the discriminants of $K,K'$.
Then both $r|_{G_{KK_0}}$ and
$r|_{G_{K'K_0}}$ are irreducible.  Otherwise
index-two Clifford theory supplies the corresponding
quadratic self-twist of $r|_{G_K}$ or
$r|_{G_{K'}}$; on inertia at $q$ the original
representation is trivial but the twisting
character is nontrivial, a contradiction.
Cyclic base change from $K$ to $KK_0$ is
cuspidal: the exceptional alternative in
Arthur--Clozel III.4.2(b) would give exactly
this prohibited self-twist.
Descend to $K_0$, base change to $K'K_0$, and
descend to $K'$ by the same argument.  All
finite-character ambiguities are corrected as
above, and all uses retain the original $\iota$.
We have obtained an actual regular algebraic
polarized cusp attached to $r|_{G_{K'}}$ for
\emph{every} field in the source's quantifier.

It remains to prove the full \texttt{WitA}, rather
than stopping at an automorphic attachment.
First note two local consequences of the actual
space.  At $p\ne\ell$, the open group
$G_F\cap\ker\bar r$ contains the commutator
subgroup of $G_{\mathbb Q}$: both $F/\mathbb Q$
and $\bar r$ are abelian.  The killing relation
therefore annihilates every $[d,u]$ with $d$ in
a decomposition group and $u$ in its inertia.
It also kills an open subgroup of inertia.
Thus the inertia image is finite and central
in the local image.  A lift of Frobenius and
this finite central inertia have dense generated
image, so the entire local image is commutative.
Its semisimplification is a sum of continuous
one-dimensional characters and its monodromy is
zero.  This proves the exact \texttt{IneAb}
condition on the original group and on every
restriction to $K'$.

At $\ell$, the flag of $X.P_y$ has the distinct
parallel weights
\[
                     e+1, e, 1, 0,\qquad e\ge2.
                     \tag{QX.25}
\]
Choose an open normal subgroup of the original
local decomposition group on which all its
diagonal characters are exactly the corresponding
powers of $\chi$; the finite-order errors in
$T.\texttt{Wts}$ disappear after shrinking.
Choose an element $u$ of its inertia on which
$\chi$ has infinite order.  Such an element
exists because the cyclotomic image of every
open local inertia subgroup is infinite.
The matrices of $u$ and $dud^{-1}$, in the
given flag basis, are upper triangular with
the same distinct diagonal entries
$\chi(u)^{e+1},\chi(u)^e,\chi(u),1$.
The equation
$r(d)r(u)=r(dud^{-1})r(d)$ forces $r(d)$
upper triangular: starting at the lower-left
corner, each entry below the diagonal is
multiplied by the nonzero difference of its
two diagonal eigenvalues, and induction makes
it zero.  Thus the ordered flag descends to
the original decomposition group.  Its diagonal
characters remain locally algebraic and parallel
with the weights (QX.25).  This proves the
actual \texttt{Flagged} condition, including
all original places; conjugating the chosen
place transports the construction.  The
restriction to any $K'$ has the same property.

Every cusp constructed above is essentially
tempered at every finite place.  The initial
$\pi_L$ is tempered by Caraiani's Theorem 1.2.
An algebraic Hecke character has absolute value
a single real power of the adelic norm, so the
inverse $\kappa$-twist preserves essential
temperedness with a single global normalization.
Strong local cyclic base change preserves and
reflects temperedness after this normalization:
in the local Langlands parameter, boundedness
on the Weil group before inserting its
$\mathrm{SL}_2$ factor is unchanged on a subgroup
of finite index.  The finitely many cosets give
the reverse implication as well.  Finite-character
twists are unitary.  This proves the assertion
through each of the descent and base-change
steps, without asserting the Ramanujan conjecture
for an unrelated unpolarized cusp form.

Here is the necessary local calculation at
$v\mid\ell$.  Write the characters in the
descended flag as
$\psi_i=\chi^{w_i}\eta_i$, where $\eta_i$
has finite inertia.  This follows from local
reciprocity and the parallel weights, or by
killing the finite inertia characters on an
open subgroup and descending back.  Continuity
makes the value of $\eta_i$ on Frobenius an
$\ell$-adic unit.  The Weil character of
$\chi$ at geometric Frobenius has value
$q^{-1}$.  Hence an equality
\[
 \operatorname{WD}(\psi_i)
   =\operatorname{WD}(\psi_j)\,|\cdot|
 \quad\Longrightarrow\quad
 w_i-w_j=1,\qquad \eta_i=\eta_j,\qquad
 \psi_i=\psi_j\chi.
 \tag{QX.26}
\]
Indeed the $\ell$-adic valuations determine the
weight difference, and the remaining Weil
characters of finite-inertia Galois characters
are faithful data.  The same implication holds
after exchanging $i,j$.  All one-dimensional
composition factors of a flag are intrinsic
as a multiset, so a different flag in
\texttt{Unl1} has the same test.

Suppose \texttt{Unl1} holds.  BLGGT's
semisimplified local compatibility identifies
the Weil constituents of the constructed cusp
with these four characters.  A nonzero
monodromy entry would give exactly a norm
ratio in (QX.26), contradicting \texttt{Unl1}.
Thus its local monodromy is zero.  Its
essential temperedness then makes all four
Weil-character absolute values equal after
one scalar normalization.  An equality
$\psi_i^m=(\psi_j\chi)^m$ with $m\ge1$
would instead give a Weil-character norm ratio
$q^{\pm m}$, a contradiction.  This proves
the first clause \texttt{Unl1}$\Rightarrow$
\texttt{PotUnl}, with every positive power
and every original decomposition group.

At a place away from $\ell$, full polarized
local compatibility identifies the Weil--Deligne
parameter with the already proved sum of
characters and zero monodromy.  At a place
above $\ell$ where \texttt{PotUnl} holds,
the preceding argument does the same.  Local
Langlands now gives a toral generic principal
series; all of its cuspidal support characters
are one-dimensional and have the same
uniformizer absolute value after the global
essential-unitary normalization.  In particular,
for one real number $t$ fixed by the global
central character, at every such place
\[
          q^{t-1/2}<|\xi(\varpi)|<q^{t+1/2}
          \quad\hbox{for every cuspidal-support character }\xi.
          \tag{QX.27}
\]
The inequalities are strict since $q>1$; the
middle value is $q^t$.  This includes the
source's \texttt{TorAt} assertion and is
exactly \texttt{WinAt}, with one $t$ for
all places.  Any other regular algebraic cusp
attached to the same original $r|_{G_{K'}}$
is isomorphic to this one by strong multiplicity
one, so the universal quantifier over such
cusps in \texttt{WitA} is also satisfied.

We have proved the actual automorphic attachment
and \texttt{WitA} over every requested imaginary
quadratic field, using the same $\iota$ throughout.
This is the literal \texttt{IsAutQA}, hence
\texttt{IsAut4Alt}, for the original $r$.
\end{proof}
% END QX9

\subsection{The exact remaining support input and the five-clause join}

Call the following assumption the \emph{rank-four ordinary-support
input on the two constructed links}.  For each original Fit and
Deep tower and each fixed $\iota$, use the coefficient threshold
and the two links constructed in QX4.  Over every further
coefficient ring, for each of these links and every actual
universal ordinary ring $B$ and weight object $\Lambda$, there
exists an ideal $J_0$ satisfying all six \texttt{HidaPropsAlt}
fields printed above and the literal $B.\texttt{P1core}\,J_0$.
This assumption is restricted to rank four and these actual
data.  It asserts existence of the ideal, including finite
torsion-free weight-module structure, classical points and
their two arithmetic endpoints, and characteristic-$\ell$
support.  It is a retained arithmetic input, not a consequence
of QX1 or a name for a theorem already proved in this chapter.

% BEGIN QX10
\begin{theorem}[QX10: the actual modular-support constructor]
\label{claim:QX10}
The stated rank-four ordinary-support input implies the entire
literal \texttt{PMExAlt}.  More explicitly, for every original
$D,T$ with $T.\texttt{Fit}$ and $T.\texttt{Deep}$ and every
fixed $\iota$, there is a coefficient threshold $c_1$ such that,
for every $c'$ containing both $c$ and $c_1$ and every actual
$X:\texttt{Space}(T.\texttt{fr}\,\varnothing)$ over $c'$, an
ideal $J\subset X.B$ satisfies all five clauses of
$T.\texttt{PotentialModularityTower}\,\iota XJ$.
\end{theorem}
\begin{proof}
Choose the common threshold of QX4; it depends on the
original tower and its character data, and may be enlarged
for $\iota$ if the retained support input requires that.
It is chosen before $c'$ and before $X$.
Fix any allowed $c'$ and any actual $X$.
For the first link, QX5 constructs an actual ordinary
\texttt{Pair} over $X.B$.  Thus the nonempty-pair
antecedent of KX9 is fulfilled, and KX9 supplies an
actual universal ring $B$.  KX5 supplies its actual
weight object $\Lambda$.
The retained input supplies $J_0$ for these data.
Perform the same construction for the second link,
obtaining $B',\Lambda',J'_0$ and the two maps
$\phi_0,\phi'_0$.  Define
\[
        J=\phi_0(J_0)X.B+\phi'_0(J'_0)X.B.
        \tag{QX.28}
\]
This is a concrete ideal in the original $X.B$;
no coefficient specialization or change of the
original space has been made.

For \texttt{fin}, use the first quotient's actual
finite weight-module field.  IX9 proves the full
algebraic finite-fibre argument: the universal
square-characteristic identities are formed
before quotienting; they make the original
characteristic coefficients integral over the
weight image.  After killing the coefficient
maximal ideal and all
$\psi_{v,i}(u)-1$ for original inertia elements,
the finite weight-module quotient is finite
dimensional over the residue field.  Its
maximal ideal is nilpotent.  Thus, for some
integer $N$,
\[
 \mathfrak m_{X.B}^{N}\subset
 J_{\mathcal L}+\mathfrak m_{c'.O}X.B+
 (\psi_{v,i}(u)-1:\ u\in I_v).
 \tag{QX.29}
\]
All hypotheses of IX9 are supplied by the actual
\texttt{RingProps.fin}, QX5's \texttt{IsPhi},
and the unchanged \texttt{Space} generation
fields.  Enlarging $J_{\mathcal L}$ to $J$
gives exactly \texttt{fin}; no injectivity of
a map after quotienting is assumed.

For \texttt{p1}, take any of its finite-field
power-series tests killing $J$.  It kills both
summands.  QX8 says that every prime below its
$X.B$-kernel contains each summand, hence their
sum.  This is the required \texttt{LocZero}.
For \texttt{seed}, QX7 applies to the same seed
point for each link.  Again every prime below
its kernel contains both summands and therefore
$J$.  Its auxiliary twisting parameter may
be chosen separately for the two proofs; QX6
returns each to its original $a=0$ image ideal.

For \texttt{pur} and \texttt{aut}, a point
killing $J$ kills its first summand.  Apply QX9
with exactly the additional hypotheses printed
in the corresponding clause.  Its proof
returns the original $r$, every place and
imaginary quadratic field, and the same fixed
$\iota$.  Thus all five fields are supplied
on this same $X,J$.
The construction worked for arbitrary $c'$
above the common threshold and arbitrary $X$;
these universal quantifiers were never
interchanged with an $X$-dependent coefficient
choice.  This is precisely \texttt{PMExAlt}.
\end{proof}
% END QX10

\paragraph{What this closes and what remains.}
QX1 supplies the literal \texttt{AdjointLiftingTheorem}.
QX2--QX3 close its fixed-weight tangent-to-support use;
QX4--QX10 close every internal step from the actual
two ordinary-support ideals to the five-clause
\texttt{PMExAlt} conclusion.  The original
\texttt{P3core} field is not used on this route.
The six-field ordinary modular quotient and its
\texttt{P1core} support remain to be constructed on
the actual rank-four residual data.  In particular
the residual-Schur and local-degree assumptions
of Allen--Newton--Thorne's residually reducible
automorphy-lifting theorem cannot be inferred from
the present Fit and Deep fields.  Neither this
chapter nor adjoint vanishing supplies those
hypotheses.  The exact Hida-control, Hecke-ideal,
arithmetic-point and characteristic-$\ell$ support
producer cut is recorded in the accompanying
frontier registry.  No unconditional \texttt{PMExAlt},
full \texttt{UnramifiedInputs}, paired-count result,
\texttt{PointHV}, RH or GRH conclusion is promoted.
